\documentclass[11pt]{article}
\usepackage[margin=1in]{geometry}
\usepackage{amsmath,amssymb,bm,graphicx,booktabs,caption,subcaption}
\usepackage[table]{xcolor}\usepackage{hyperref}\hypersetup{colorlinks=true,linkcolor=blue!50!black}
\graphicspath{{figs/}}\captionsetup{font=small}
\title{\textbf{A recurrent vision transformer shows signatures of primate visual attention:\\
a full-figure reproduction on a nine-stimulus grid with a multiplicative-gating variant}}
\author{Reproduction of Morgan, Albanna \& Herman with the conv-frontend, $d_{\mathrm{mem}}{=}128$ \emph{element-wise multiplicative-gating} JEPA variant on the 9-stimulus \texttt{vda9} task}
\date{Model checkpoint iteration 2{,}999 (task: 9-stimulus $3\times3$ grid \texttt{vda9}, $75\times75$)}
\begin{document}\maketitle
\begin{abstract}\noindent
We reproduce \emph{every} figure of Morgan, Albanna \& Herman on a harder, nine-stimulus version of the
cued orientation-change task---a $3\times3$ Gabor grid with cues at opposite corners ($S_1$ top-left,
$S_9$ bottom-right)---using a convolutional-frontend, $d_{\mathrm{mem}}{=}128$ multiplicative-gating JEPA
variant ($X'{=}\gamma(H)\odot X+\beta(H)$) trained from pixels by reinforcement learning. The variant produces orderly sigmoidal
psychometric and chronometric functions, a clear cued$>$uncued cueing benefit, interpretable recurrent
attention dynamics over the $3\times3$ field, decodable change location/occurrence in the mnemonic
percept, a competitive actor-logit decision axis, a distributional value jump at the change, and the
signal-detection criterion/sensitivity signatures of attentional ``stimulation.'' Unlike the compact
four-stimulus variants---which ignore displayed validity and (for the multiplicative variant) have inert
hard clamps---this nine-stimulus multiplicative model \emph{recruits validity} on three fronts: a mild
cued-threshold decrease with validity, a cued$>$uncued benefit that grows with validity, and a
validity-graded re-allocation of attention onto the cued corner ($\alpha_1$ at the change frame
$0.15\to0.48$ from $0\%$ to $100\%$); and its attention is a genuine causal lever (a clamp ``dual
effect'' and a parametric threshold shift). We report all figures at an early ($2{,}999$-iteration)
checkpoint. \emph{This is a mid-training snapshot; the qualitative signatures are already present and
stable, and where behaviour is still maturing we say so.}
\end{abstract}

\section{Introduction}
The nine-stimulus grid raises set size from four to nine, distributing the same cue/change structure
over a $3\times3$ field. We ask whether the primate-like attention signatures of the original---and of
our four-stimulus reproductions---survive this harder task, and in particular whether the larger field
recruits displayed cue validity more than the four-stimulus variants did (which ignored it entirely).

\section{Task and model}
Each 7-step trial (Fig.~\ref{fig:task}) presents a $75\times75$ image: blank, a corner cue at $S_1$ or
$S_9$ ($t{=}1$), blank, then nine Gabor patches from $t{=}3$; on half of trials one patch's orientation
steps by $\Delta$ at $t{=}5$. The variant (Fig.~\ref{fig:model}) parses each frame into nine patches,
encodes them with a shared SE-ResNet, and feeds the tokens to a recurrent transformer whose nine visual
tokens are affinely gated by the memory ($X'{=}\gamma(H)\odot X+\beta(H)$) before self-attending; a
spatial xLSTM holds the working memory read out to a distributional actor--critic and JEPA head.

\begin{figure}[h]\centering
\includegraphics[width=\textwidth]{fig1A_task_timeline.png}\\[4pt]
\includegraphics[width=0.6\textwidth]{fig1B_cue_configs.png}
\caption{\textbf{Fig 1 --- Task / environment (9-stim).} Top: the 7-step timeline with change/no-change
branch. Bottom: cue configurations (validity ring $\times$ corner) with the empirical $3\times3$
change-location probability our environment produces.}\label{fig:task}\end{figure}

\begin{figure}[h]\centering\includegraphics[width=\textwidth]{fig2_schematic_affine_ew.png}
\caption{\textbf{Fig 2 --- Model schematic (9-stim multiplicative-gating).} Nine-patch field $\to$ conv
SE-ResNet $\to$ recurrent-ViT self-attention with memory gating ($X'{=}\gamma(H)\odot X+\beta(H)$) $\to$
spatial xLSTM $\to$ actor--critic (RL) and JEPA head.}\label{fig:model}\end{figure}

\section{Behaviour: psychometrics and chronometrics}
The variant produces orderly sigmoidal psychometrics and decreasing chronometrics
(Fig.~\ref{fig:psy}). \emph{Validity.} The fitted threshold for cued changes shows a \emph{mild}
validity-graded decrease ($15.8/16.1/14.4/13.9/14.2^\circ$ for $0/25/50/75/100\%$)---a weak but present
threshold effect, stronger than in either four-stimulus variant. And---unlike the four-stimulus
variants---the \emph{cued$>$uncued benefit grows with validity}: at a $0\%$ (uninformative) cue the
cued and uncued curves nearly coincide, whereas at $100\%$ validity the cued curve is markedly left of
the uncued (Fig.~\ref{fig:psy}B vs.\ C). The nine-stimulus model therefore recruits displayed validity
\emph{partially}, in the spatial split rather than in the cued threshold.

\begin{figure}[h]\centering\includegraphics[width=\textwidth]{fig3_affine_ew.png}
\caption{\textbf{Fig 3 --- psychometrics and chronometrics.} A: cued response vs.\ $\Delta$ at five
validities (incl.\ $0\%$). B: $0\%$ cued vs.\ uncued. C: $100\%$ cued vs.\ uncued (benefit larger).
D--F: reaction time. Points $n{=}300$; lines are 4-parameter logistic fits.}\label{fig:psy}\end{figure}
\begin{figure}[h]\centering\includegraphics[width=0.7\textwidth]{fig3_affine_ew_params.png}
\caption{\textbf{Fitted psychometric parameters} (cued) vs.\ validity.}\label{fig:psyp}\end{figure}

\section{Attention allocation dynamics}
Averaged recurrent self-attention over the $3\times3$ field (Fig.~\ref{fig:attn}) is flat at onset
($\approx\!1/9$ per patch) and rises onto the cued corner at the cue frame; because attention here enters
through a multiplicative gate rather than a competitive key/value, the cue peak is more modest than in
the cross-attention variant ($\alpha_1{\approx}0.40$ vs.\ $0.89$). Crucially---and \emph{unlike} the
four-stimulus variants---the \emph{re-allocation} of attention onto the cued corner during the stimulus
epoch is graded by displayed validity, and strongly so: at the change frame $\alpha_1$ grows
monotonically from $\approx\!0.15$ at $0\%$ validity to $\approx\!0.48$ at $100\%$. Thus the
nine-stimulus model's \emph{attention} is validity-modulated even where its cued psychometric threshold
is not, mirroring the behavioural finding that the cued$>$uncued benefit widens with validity. Panel B
shows $\alpha_1$($S_1$)/$\alpha_9$($S_9$) vs.\ time; panel C shows $\alpha_1$ at the change frame vs.\
$\Delta$ for cued and uncued changes.

\begin{figure}[h]\centering
\includegraphics[width=\textwidth]{fig4A_affine_ew.png}\\[4pt]
\includegraphics[width=0.47\textwidth]{fig4B_affine_ew.png}\hfill
\includegraphics[width=0.47\textwidth]{fig4C_affine_ew.png}
\caption{\textbf{Fig 4 --- intrinsic attention dynamics (9-stim).} Top: per-timestep $3\times3$
attention heatmaps by validity (no-change cue-$S_1$). Bottom left: $\alpha_1$/$\alpha_9$ vs.\ time.
Bottom right: $\alpha_1$ at change vs.\ $\Delta$, uncued/cued.}\label{fig:attn}\end{figure}

\section{Feedback mechanism}
Our variant instantiates the \emph{multiplicative} feedback, gating the nine visual tokens element-wise
by the memory ($X'{=}\gamma(H)\odot X+\beta(H)$) before self-attention---the mechanism the original
selected as its focus. Fig.~\ref{fig:fb} shows the circuit alongside the four behavioural panels of the
original's feedback figures.

\begin{figure}[h]\centering\includegraphics[width=\textwidth]{fig567_affine_ew.png}
\caption{\textbf{Figs 5--7 analog --- feedback circuit \& behaviour (9-stim).} Left: multiplicative-gating
circuit. A: cued response by validity. B: enhance/suppress attention on the cued change. C: cued vs.\
uncued at $100\%$. D: attention on $S_1$(cued) vs.\ $S_9$(change) at the change frame.}\label{fig:fb}\end{figure}

\section{Causal attention manipulations on behaviour}
Clamping attention at the change frame (Fig.~\ref{fig:manip}) reproduces the ``dual effect'': for an
uncued change ($S_9$, $100\%$, $\Delta{=}18^\circ$), directing attention onto the change ($\alpha_9{=}1$)
raises detection ($0.69\to0.89$) and \emph{suppressing} the cued location ($\alpha_1{=}0$) \emph{also}
raises it ($0.69\to0.78$). Notably---and \emph{in contrast to the four-stimulus multiplicative
variant, whose hard clamps were inert}---here the clamps are a genuine causal lever, and the parametric
$\alpha_1{\leftrightarrow}\alpha_9$ sweep grades the psychometric threshold ($13.2^\circ\!\to\!16.0^\circ$
as $\alpha_1$ falls $1.0\to0.1$; Fig.~\ref{fig:manipp}). On the harder nine-stimulus task the
multiplicative variant thus learns a more behaviourally load-bearing attention than its four-stimulus
counterpart.

\begin{figure}[h]\centering\includegraphics[width=\textwidth]{fig5AF_affine_ew.png}
\caption{\textbf{Attention clamps at $t_{\mathrm{change}}$ (9-stim).} Response (A--C) and RT (D--F) for
clamps on $S_1$/$S_9$.}\label{fig:manip}\end{figure}
\begin{figure}[h]\centering\includegraphics[width=\textwidth]{fig5GK_affine_ew.png}
\caption{\textbf{Parametric manipulation (9-stim).} Shifting attention $S_1{\leftrightarrow}S_9$ grades
the psychometric function; fitted parameters vs.\ $\alpha_1$.}\label{fig:manipp}\end{figure}

\section{Decoding the mnemonic percept}
Decoders trained on natural states recover change \emph{location} (10-way over the $3\times3$ grid $+$
no-change) from the full percept $H$ (Fig.~\ref{fig:dec8}) and change \emph{occurrence} from a single
patch $h_1$ (Fig.~\ref{fig:dec9}); the actor's first activation retains occurrence
(Fig.~\ref{fig:dec10}), and graded inhibition of $\alpha_1$ is examined in Fig.~\ref{fig:dec11}.

\begin{figure}[h]\centering\includegraphics[width=\textwidth]{fig8_affine_ew.png}
\caption{\textbf{Fig 8 --- decode change LOCATION (10-way) from $H$}, $t{=}5$/$6$, under
\{no-mod, spread, $\alpha_1{=}0$, $\alpha_1{=}1$\}.}\label{fig:dec8}\end{figure}
\begin{figure}[h]\centering\includegraphics[width=0.85\textwidth]{fig9_affine_ew.png}
\caption{\textbf{Fig 9 --- decode change OCCURRENCE from single patch $h_1$}, $t{=}5$/$6$.}\label{fig:dec9}\end{figure}
\begin{figure}[h]\centering\includegraphics[width=\textwidth]{fig10_affine_ew.png}
\caption{\textbf{Fig 10 --- decode from actor first activation} (location/occurrence).}\label{fig:dec10}\end{figure}
\begin{figure}[h]\centering\includegraphics[width=\textwidth]{fig11_affine_ew.png}
\caption{\textbf{Fig 11 --- decode $S_1$-change vs.\ no-change from actor activation}, graded
$\alpha_1$.}\label{fig:dec11}\end{figure}

\section{Actor-logit geometry}
The actor logits at the change frame form a competitive Declare/Wait axis graded by $\Delta$
(Fig.~\ref{fig:log12}), shifted by attention modulation on $S_1$ (Fig.~\ref{fig:log13}).

\begin{figure}[h]\centering\includegraphics[width=0.62\textwidth]{fig12_affine_ew.png}
\caption{\textbf{Fig 12 --- actor logit geometry (9-stim).} A: natural. B--E: changes at $S_1$/$S_3$/
$S_5$/$S_9$. Coloured by $\Delta$ (left) and density (right).}\label{fig:log12}\end{figure}
\begin{figure}[h]\centering\includegraphics[width=0.62\textwidth]{fig13_affine_ew.png}
\caption{\textbf{Fig 13 --- actor logits under $S_1$ attention modulation.}}\label{fig:log13}\end{figure}

\section{Value, TD error and entropy}
The distributional critic's value grows with $\Delta$ (Fig.~\ref{fig:td14}), and entropy/criterion vs.\
attention are shown in Fig.~\ref{fig:ent15}.

\begin{figure}[h]\centering\includegraphics[width=\textwidth]{fig14_affine_ew.png}
\caption{\textbf{Fig 14 --- TD error \& value (9-stim).}}\label{fig:td14}\end{figure}
\begin{figure}[h]\centering\includegraphics[width=0.95\textwidth]{fig15_affine_ew.png}
\caption{\textbf{Fig 15 --- entropy \& criterion--value coupling (9-stim).}}\label{fig:ent15}\end{figure}

\section{Signal detection: criterion \& sensitivity (``stimulation'' experiments)}
Biasing attention on $S_1$ parametrically (at cue-time or change-time) modulates the signal-detection
criterion and sensitivity (Fig.~\ref{fig:sdt16}): attention toward the change lowers criterion and
preserves sensitivity, whereas attention directed away reduces sensitivity.

\begin{figure}[h]\centering\includegraphics[width=\textwidth]{fig16_affine_ew.png}
\caption{\textbf{Fig 16 --- SDT criterion \& sensitivity vs.\ attention bias on $S_1$ (9-stim).} Rows:
$c$, $d'$. Columns: manipulate at cue-time / change-time. Lines: \{Cue $25/100\%$\}$\times$\{change
$S_1$/$S_9$\}.}\label{fig:sdt16}\end{figure}

\section{Supervised vs.\ RL}
We reproduce the RL column across the four behavioural rows; the supervised columns are out-of-scope
baselines not trained for this RL-only variant (Fig.~\ref{fig:sup17}).

\begin{figure}[h]\centering\includegraphics[width=0.9\textwidth]{fig17_affine_ew.png}
\caption{\textbf{Fig 17 analog --- supervised vs.\ RL (9-stim).}}\label{fig:sup17}\end{figure}

\section{Discussion}
On the harder nine-stimulus grid the multiplicative-gating variant reproduces the full battery of
primate-like signatures and, strikingly, recruits displayed cue validity far more than any of the
four-stimulus variants: a mild cued-threshold decrease, a cued$>$uncued benefit that grows with
validity, and a validity-graded re-allocation of attention onto the cued corner. Its attention is also a
genuine causal lever here---a clamp ``dual effect'' and a parametric threshold shift---whereas the
\emph{same} multiplicative mechanism on the four-stimulus task had inert hard clamps. Set size and task
difficulty thus appear to push this architecture toward using attention (and validity) more, not less.
Because this is an early ($2{,}999$-iteration) checkpoint, we present these as robust qualitative
observations rather than converged quantitative effects; continued training is expected to sharpen them.

\section{Methods (variant)}
As in the four-stimulus reproduction but with a $3\times3=9$-patch field ($75\times75$ image, corner
cues $S_1$/$S_9$). Front-end: shared SE-ResNet to $\hat o_i\in\mathbb R^{128}$. Multiplicative gating
$X'{=}\gamma(H)\odot X+\beta(H)$ then self-attention over the nine tokens, $Z=X'+AV$. Memory: per-patch
spatial xLSTM, $d_{\mathrm{mem}}{=}128$. Training:
distributional actor--critic on sparse reward $+$ V-JEPA. Analyses use $n{=}300$ trials/point;
psychometrics fit a 4-parameter logistic. Decoders are 10-way (location) or 2-way (occurrence), trained
on natural states. SDT uses hit/false-alarm rates at $\Delta{=}18^\circ$.

\end{document}
